% =====================================================================
%  Section 4 -- Markov model of exposure (owner: THEORY).
%  Proofs in appendix/A-proofs-estimands.tex. Table values: invented matrices,
%  computed deterministically (numpy); see the THEORY hand-off note.
% =====================================================================
\providecommand{\osc}{\operatorname{osc}}

\section{A Markov model of exposure}\label{sec:markov}

\Cref{def:amplification} is model-free. A model is needed for three tasks: to extrapolate from the audited session length to longer sessions and to the long run, to compare policies whose sessions start differently, and to understand why a small change in a recommender can produce a large change in exposure. We model a session as a Markov chain on content states. The tools are classical: the fundamental matrix of \citet{kemeny1960}, the group inverse of \citet{meyer1975}, and the perturbation theory of finite chains \citep{schweitzer1968,meyer1980,seneta1988,cho2000,cho2001}. \citet{ribeiro2020} studied, on graphs of recommendations, how readily one community of channels can be reached from another. We use the chain to compute exposure functionals and their sensitivity. \Cref{prop:markov} collects the exposure functionals, \cref{thm:perturb} bounds their sensitivity, and \cref{prop:markov-mle} gives their estimation from audit sessions. \Cref{tab:markov-illus} illustrates the first two with invented matrices. Parts (a) and (b) of \cref{thm:perturb} restate classical results for exposure functionals. The closed form in part (c) is the personalised PageRank vector with teleportation to a single state \citep{page1999,jeh2003}. We have not found the hitting-time form of the derivative in part (c), or parts (d) and (e), stated in this form, although each follows in a few lines from the classical theory.

\subsection{Sessions as Markov chains}\label{sec:markov-model}

Let $\mathcal X=\{1,\dots,m\}$ be a finite set of content states, partitioned into neutral states $\mathcal X_N$, borderline or adjacent states $\mathcal X_B$, and antisemitic states $\mathcal A$. The coarsest useful choice has one state of each kind, and finer choices cross severity with topic. Only the target set $\mathcal A$ matters for the theory. A session is a sequence of states $X_1,\dots,X_T$, the states of the items $c_1,\dots,c_T$, so that $\ind{L_{c_t}=1}=\ind{X_t\in\mathcal A}$.

\begin{assumption}[Markov sessions]\label{asm:markov}
Given the policy $\pi$ and the persona $d$, the states of a session form a time-homogeneous Markov chain with initial distribution $\alpha_\pi(d)$ and transition matrix $P_\pi(d)$, where $P_\pi(d)_{xy}$ is the probability that an item in state $x$ is followed by an item in state $y$.
\end{assumption}

The transition matrix combines supply and demand: the policy determines what is offered after an item in state $x$, and the persona determines what is taken. We suppress $(\pi,d)$ when they are fixed and write $P_{\mathrm{alg}}$ and $P_{\mathrm{ref}}$ when comparing policies at the same persona. \Cref{asm:markov} is an approximation. It can be made to hold more nearly by enlarging the state, for example to the last $k$ states or to the pair (state, position in session). The order of the chain can be tested by likelihood-ratio tests \citep{anderson1957}. With $\bm1$ the vector of ones and $\bm1_{\mathcal A}$ the indicator of $\mathcal A$, the expected exposure \eqref{eq:est-exposure} of a $T$-item session is
\begin{equation}\label{eq:mk-eT}
  e_T(P,\alpha)=\E\Bigl\{\frac1T\sum_{t=1}^{T}\ind{X_t\in\mathcal A}\Bigr\}=\frac1T\sum_{t=0}^{T-1}\alpha^{\top}P^{t}\bm1_{\mathcal A},
\end{equation}
and $\mu_\pi(d)=e_T\{P_\pi(d),\alpha_\pi(d)\}$.

For a stochastic matrix $P$ with stationary distribution $\varpi$ (a column vector), let
\[
  Z=(I-P+\bm1\varpi^{\top})^{-1},\qquad H=Z-\bm1\varpi^{\top},\qquad \bm g_{\mathcal A}=H\bm1_{\mathcal A}.
\]
$Z$ is the fundamental matrix of \citet{kemeny1960}, and $H$ is the group inverse $(I-P)^{\#}$ of \citet{meyer1975}, also called the deviation matrix. We call $\bm g_{\mathcal A}$ the \emph{exposure potential}. It solves the Poisson equation $(I-P)\bm g_{\mathcal A}=\bm1_{\mathcal A}-\varpi(\mathcal A)\bm1$, and for aperiodic $P$, $g_{\mathcal A}(y)=\sum_{t\ge0}\{\Pr_y(X_{1+t}\in\mathcal A)-\varpi(\mathcal A)\}$ is the excess antisemitic exposure, summed over the rest of the session, of a session that passes through state $y$. For a set $\mathcal S\subseteq\mathcal X$ and a chain started at $X_1=x$, let $\tau_{\mathcal S}=\inf\{n\ge0:X_{1+n}\in\mathcal S\}$ be the number of further items until the session first reaches $\mathcal S$, let $M_{x\mathcal S}=\E_x\tau_{\mathcal S}$, and write $M_{xy}=M_{x\{y\}}$, so that $M_{yy}=0$. Finally, $\|B\|_\infty=\max_x\sum_y|B_{xy}|$ is the maximum absolute row sum, $\tau_1(B)=\frac12\max_{x,y}\sum_z|B_{xz}-B_{yz}|$ is the ergodicity coefficient of \citet{seneta2006}, and $\osc(\bm v)=\max_xv_x-\min_xv_x$.

\subsection{Exposure functionals}\label{sec:markov-functionals}

A stochastic matrix is \emph{unichain} if it has exactly one closed communicating class. It then has a unique stationary distribution, and states outside the class are transient.

\begin{proposition}[Exposure functionals]\label{prop:markov}
Let $P$ be a stochastic matrix. In parts (a)--(c), $P$ is unichain with stationary distribution $\varpi$, and $Z$, $H$ and $\bm g_{\mathcal A}$ are as above. Part (d) makes no such assumption.
\begin{enumerate}[(a)]
  \item \emph{Long run.} $Z$ exists, $H$ is the group inverse of $I-P$, and $\varpi^{\top}=\bm1^{\top}(I-P+\bm1\bm1^{\top})^{-1}$. For every initial distribution, the session exposure $E=T^{-1}\sum_{t=1}^{T}\ind{X_t\in\mathcal A}$ converges to $\varpi(\mathcal A)$ almost surely, and $\sqrt T\{E-\varpi(\mathcal A)\}$ converges in distribution to $\mathrm N(0,\sigma^2_{\mathcal A})$, where, with $\tilde{\bm f}=\bm1_{\mathcal A}-\varpi(\mathcal A)\bm1$,
  \[
    \sigma^2_{\mathcal A}=2\sum_x\varpi_x\tilde f_x(H\tilde{\bm f})_x-\sum_x\varpi_x\tilde f_x^{\,2}
    =\sum_x\varpi_x\Var_{y\sim P(x,\cdot)}\{g_{\mathcal A}(y)\}.
  \]
  \item \emph{Finite horizon.} For every $\alpha$ and $T\ge1$,
  \[
    e_T(P,\alpha)=\varpi(\mathcal A)+\frac1T\,\alpha^{\top}(I-P^{T})\bm g_{\mathcal A}
    =\varpi(\mathcal A)+\frac1T\,\alpha^{\top}(I-P^{T})Z\bm1_{\mathcal A}.
  \]
  Hence $e_T(P,\alpha)=\varpi(\mathcal A)+b(\alpha)/T+r_T/T$, where $b(\alpha)=\alpha^{\top}\bm g_{\mathcal A}=(\alpha-\varpi)^{\top}Z\bm1_{\mathcal A}$ and $|r_T|\le\frac12\osc(\bm g_{\mathcal A})\,\|\alpha^{\top}P^{T}-\varpi^{\top}\|_1$, which tends to zero if $P$ is aperiodic.
  \item \emph{Hitting.} Let $\mathcal O=\mathcal X\setminus\mathcal A$. The probability that a $T$-item session contains no antisemitic item is $\alpha_{\mathcal O}^{\top}P_{\mathcal O\mathcal O}^{T-1}\bm1$. If $\mathcal A$ is accessible from every state, then $I-P_{\mathcal O\mathcal O}$ is invertible and $M_{x\mathcal A}=\{(I-P_{\mathcal O\mathcal O})^{-1}\bm1\}_x$ for $x\in\mathcal O$. If $P$ is irreducible, then $M_{xy}=(H_{yy}-H_{xy})/\varpi_y$ for all $x,y$, and
  \[
    K(P)=\sum_y\varpi_yM_{xy}=\operatorname{tr}H=\sum_{k=2}^{m}\frac{1}{1-\lambda_k}
  \]
  does not depend on $x$. Here $\lambda_2,\dots,\lambda_m$ are the eigenvalues of $P$ other than the simple eigenvalue $1$, counted with algebraic multiplicity.
  \item \emph{Rabbit holes.} Let $P$ be any stochastic matrix, with closed classes $C_1,\dots,C_r$ ($r\ge1$) and transient set $\mathcal T$, and let $\varpi^{(k)}$ be the stationary distribution of $P$ restricted to $C_k$. Then $N=(I-P_{\mathcal T\mathcal T})^{-1}$ exists. From $x\in\mathcal T$ the session enters $C_k$ with probability $B_{xk}=(NP_{\mathcal TC_k}\bm1)_x$, after an expected $(N\bm1)_x$ items, and for every $\alpha$,
  \[
    \lim_{T\to\infty}e_T(P,\alpha)=\sum_{k=1}^{r}\bigl\{\alpha(C_k)+\alpha_{\mathcal T}^{\top}NP_{\mathcal TC_k}\bm1\bigr\}\,\varpi^{(k)}(\mathcal A).
  \]
\end{enumerate}
\end{proposition}

Parts (a), (c) and (d) are classical \citep[Chs.~3--4]{kemeny1960}, and the group-inverse forms are those of \citet{meyer1975}. Kemeny and Snell give the full limiting covariance matrix of the occupation counts of the states \citep[Thm.~4.6.1]{kemeny1960}, of which $\sigma^2_{\mathcal A}$ is a quadratic form. The second expression for $\sigma^2_{\mathcal A}$, a martingale form, is the one used for estimation in \cref{sec:markov-estimation}. $K(P)$ is Kemeny's constant. With the convention $M_{yy}=0$ it is the expected number of items needed to reach a state drawn from $\varpi$, from any starting state. Part (b) is exact for every $T$. The $1/T$ term $b(\alpha)$ is the transient excess of a session that starts at $\alpha$ rather than in equilibrium. It is negative when sessions start in states of below-average exposure potential, as neutral starting states typically are, and short sessions then understate long-run exposure.

Part (d) gives the formal meaning of a \emph{rabbit hole}: a closed class $C$ that contains an antisemitic state and no neutral state, $C\cap\mathcal A\neq\emptyset$ and $C\cap\mathcal X_N=\emptyset$. Once a session enters such a class it never returns to neutral content, and its long-run exposure is $\varpi^{(C)}(\mathcal A)$, however rarely the class is entered. A strict rabbit hole is a modelling idealisation, because a reference policy that occasionally shows neutral content breaks it. The relevant case in practice is a \emph{nearly} closed class, which part (d) of \cref{thm:perturb} addresses.

\subsection{Sensitivity: small nudges, large exposure}\label{sec:markov-perturb}

Let the algorithm differ from the reference by a change $P_{\mathrm{alg}}-P_{\mathrm{ref}}$ in recommendation probabilities. Along the path $P(\varepsilon)=P_{\mathrm{ref}}+\varepsilon\Delta$, with $\Delta\bm1=\bm0$, the parameter $\varepsilon$ measures the size of the nudge and $\Delta$ its direction. Any algorithm can be reached this way, with $\Delta=(P_{\mathrm{alg}}-P_{\mathrm{ref}})/\varepsilon$. The question is how much long-run and finite-horizon exposure can change per unit of nudge. Define the \emph{sensitivity coefficient} of $\mathcal A$ under a chain $P$ as
\[
  \kappa_{\mathcal A}(P)=\tfrac12\osc(\bm g_{\mathcal A}),\qquad \kappa_y=\kappa_{\{y\}} .
\]

\begin{theorem}[Sensitivity of exposure]\label{thm:perturb}
Let $P_{\mathrm{ref}}$ be irreducible, with stationary distribution $\varpi$ and with $Z$, $H$, $\bm g_{\mathcal A}$, $\kappa_{\mathcal A}$ and $M$ computed from $P_{\mathrm{ref}}$. Let $P_{\mathrm{alg}}$ be unichain with stationary distribution $\varpi_{\mathrm{alg}}$.
\begin{enumerate}[(a)]
  \item \emph{Exact change and derivative.} $\varpi_{\mathrm{alg}}^{\top}-\varpi^{\top}=\varpi_{\mathrm{alg}}^{\top}(P_{\mathrm{alg}}-P_{\mathrm{ref}})H$. The matrix $I-(P_{\mathrm{alg}}-P_{\mathrm{ref}})H$ is invertible, and $\varpi_{\mathrm{alg}}^{\top}=\varpi^{\top}\{I-(P_{\mathrm{alg}}-P_{\mathrm{ref}})H\}^{-1}$. Along a path $P(\varepsilon)$ of unichain matrices, $\varpi(\varepsilon)$ is a rational function of $\varepsilon$ with
  \[
    \frac{d\varpi(\varepsilon)^{\top}}{d\varepsilon}=\varpi(\varepsilon)^{\top}\Delta H(\varepsilon)=\varpi(\varepsilon)^{\top}\Delta Z(\varepsilon),\qquad
    \frac{d\varpi_{\mathcal A}}{d\varepsilon}\Big|_{\varepsilon=0}=\sum_x\varpi_x\sum_y\Delta_{xy}\,g_{\mathcal A}(y).
  \]
  \item \emph{Global bounds.} $|\varpi_{\mathrm{alg}}(\mathcal A)-\varpi(\mathcal A)|\le\kappa_{\mathcal A}\,\|P_{\mathrm{alg}}-P_{\mathrm{ref}}\|_\infty$ and $\|\varpi_{\mathrm{alg}}-\varpi\|_1\le\tau_1(H)\,\|P_{\mathrm{alg}}-P_{\mathrm{ref}}\|_\infty$. Moreover
  \[
    \kappa_{\mathcal A}=\tfrac12\osc_x\Bigl(\sum_{y\in\mathcal A}\varpi_yM_{xy}\Bigr)\le\tfrac12\min\{\tau_1(H),K(P_{\mathrm{ref}})\},
    \qquad \kappa_y=\tfrac12\varpi_y\max_xM_{xy}.
  \]
  Over formal directions, with $\Delta\bm1=\bm0$ and $\|\Delta\|_\infty\le1$, the supremum of $|d\varpi_{\mathcal A}/d\varepsilon|_{\varepsilon=0}$ equals $\kappa_{\mathcal A}$. It is attained by moving mass $\frac12$ in every row from a minimiser to a maximiser of $\bm g_{\mathcal A}$. Over directions that keep $P_{\mathrm{ref}}+\varepsilon\Delta$ stochastic for small $\varepsilon>0$, the largest first-order increase is $\frac12\sum_x\varpi_x\{\max_yg_{\mathcal A}(y)-\min_{y:P_{\mathrm{ref}}(x,y)>0}g_{\mathcal A}(y)\}$. This equals $\kappa_{\mathcal A}$ when every row of $P_{\mathrm{ref}}$ charges a minimiser of $\bm g_{\mathcal A}$, for instance when $P_{\mathrm{ref}}>0$, and is smaller otherwise.
  \item \emph{Nudges toward an antisemitic state.} Let $j\in\mathcal A$ and $P(\varepsilon)=(1-\varepsilon)P_{\mathrm{ref}}+\varepsilon\bm1\bm e_j^{\top}$: with probability $\varepsilon$, the next item is replaced by an item in state $j$. Then $\varpi(\varepsilon)^{\top}=\varepsilon\,\bm e_j^{\top}\{I-(1-\varepsilon)P_{\mathrm{ref}}\}^{-1}$ for $0<\varepsilon\le1$, and
  \[
    \frac{d\log\varpi_{\mathcal A}(\varepsilon)}{d\varepsilon}\Big|_{\varepsilon=0}=\frac{g_{\mathcal A}(j)}{\varpi(\mathcal A)}
    =\frac{1}{\varpi(\mathcal A)}\sum_{y\in\mathcal A}\varpi_y\bigl(\E_\varpi[\tau_y]-M_{jy}\bigr),
  \]
  where $\E_\varpi[\tau_y]=\sum_x\varpi_xM_{xy}$ is the mean number of items needed to reach $y$ from equilibrium. If $\mathcal A=\{j\}$, the derivative equals $\E_\varpi[\tau_j]$. For a memoryless stream, $P_{\mathrm{ref}}=\bm1\varpi^{\top}$, it equals $\{1-\varpi(\mathcal A)\}/\varpi(\mathcal A)$.
  \item \emph{Stickiness.} Let $j\in\mathcal X$, let $S\subseteq\mathcal X\setminus\{j\}$ be nonempty, and let $q(S)=\max_{x\in S}P_{\mathrm{ref}}(x,\mathcal X\setminus S)$ be the largest one-step probability of leaving $S$. Then $M_{xj}\ge1/q(S)$ for every $x\in S$, so that
  \[
    \kappa_j\ge\frac{\varpi_j}{2q(S)},\qquad \E_\varpi[\tau_j]\ge\frac{\varpi(S)}{q(S)}.
  \]
  \item \emph{Finite horizon.} For any stochastic matrices $P_{\mathrm{ref}}$ and $P_{\mathrm{alg}}$, every $\alpha$ and $T\ge2$,
  \[
    |e_T(P_{\mathrm{alg}},\alpha)-e_T(P_{\mathrm{ref}},\alpha)|
    \le\frac{\|P_{\mathrm{alg}}-P_{\mathrm{ref}}\|_\infty}{2T}\sum_{r=0}^{T-2}(T-1-r)\,\tau_1(P_{\mathrm{ref}}^{\,r})
    \le\frac{T-1}{4}\,\|P_{\mathrm{alg}}-P_{\mathrm{ref}}\|_\infty,
  \]
  and along $P(\varepsilon)$, $de_T/d\varepsilon|_{\varepsilon=0}=T^{-1}\sum_{t=1}^{T-1}\sum_{s=0}^{t-1}\alpha^{\top}P_{\mathrm{ref}}^{\,s}\Delta P_{\mathrm{ref}}^{\,t-1-s}\bm1_{\mathcal A}$.
\end{enumerate}
\end{theorem}

The identity in (a) is due to \citet{schweitzer1968}, who wrote it with $Z$, and to \citet{meyer1975,meyer1980} in group-inverse form. It is exact, not a first-order approximation, so the long-run exposure under the algorithm can be computed from the reference chain and the change in transition probabilities alone. The derivative has a direct reading. The algorithm moves recommendation mass $\Delta_{xy}$ from state $x$, which is visited with frequency $\varpi_x$, to state $y$, and every unit moved to $y$ is worth $g_{\mathcal A}(y)$ units of long-run exposure. Redirecting mass matters in proportion to the difference in exposure potential between the states that gain and the states that lose.

Part (b) turns this into a condition number. The $\ell_1$ bound is \citeauthor{seneta1988}'s (\citeyear{seneta1988}), and the single-state form of the first bound is the mean-first-passage bound of \citet{cho2000}. \Citet{cho2001} compare these and other condition numbers. Kemeny's constant, read as an expected time to mixing, is applied to perturbed chains by \citet{hunter2006}, and perturbation bounds in terms of Kemeny's constant and mean first passage times are also discussed by \citet{ortner2024}. For exposure, the useful form is the middle expression: $\kappa_{\mathcal A}$ is large only if some state is far, in expected items, from the antisemitic states, and it is at most half of Kemeny's constant. Part (d) gives a partial converse. If some region $S$ that excludes an antisemitic state $j$ is hard to leave, and $j$ still carries non-negligible long-run mass, then $\kappa_j$ is large. A non-negligible $\varpi_j$ despite rare entry means that $j$, once reached, is hard to leave as well. The chain then has two nearly closed regions joined by a bottleneck, the nearly decomposable structure of \citet{simon1961}. In the limit of a closed class, the bottleneck becomes a rabbit hole (\cref{prop:markov}(d)) and the stationary distribution jumps.

Part (c) answers the question in the section title. When a fraction $\varepsilon$ of recommendations is redirected to an antisemitic state $j$, long-run exposure rises by the relative amount $\varepsilon\,g_{\mathcal A}(j)/\varpi(\mathcal A)$ to first order. For aperiodic $P_{\mathrm{ref}}$, $g_{\mathcal A}(j)$ is the excess antisemitic exposure that one injected item in state $j$ generates over the rest of the session. With a single antisemitic state the multiplier is $\E_\varpi[\tau_j]$, the mean time to reach $j$ from equilibrium. With several, it averages $\E_\varpi[\tau_y]-M_{jy}$ over $y\in\mathcal A$ with weights $\varpi_y/\varpi(\mathcal A)$: for each antisemitic state $y$, an injected item in state $j$ saves the difference between the time needed to reach $y$ from equilibrium and from $j$. The multiplier can therefore be much smaller than $\E_\varpi[\tau_j]$ when $j$ does not lead quickly to the other antisemitic states. For a memoryless stream the multiplier is $\{1-\varpi(\mathcal A)\}/\varpi(\mathcal A)$, the trivial effect of adding $\varepsilon$ to the rate. The ratio $g_{\mathcal A}(j)/\{1-\varpi(\mathcal A)\}$ therefore measures how much stickiness magnifies a nudge. For $\mathcal A=\{j\}$, part (d) shows that it is large when the reference keeps users in neutral territory and the antisemitic state, once reached, keeps them in its region.

\begin{example}[Two regions]\label{ex:mk-two-region}
Consider a two-state chain on $\{N,R\}$, where $R$ collects borderline and antisemitic content, with entry probability $a=P(N,R)$ and exit probability $\delta=P(R,N)$. Then $\varpi_R=a/(a+\delta)$, $K=1/(a+\delta)$, $\kappa_R=1/\{2(a+\delta)\}$ and $d\varpi_R/da=\delta/(a+\delta)^2=(1-\varpi_R)^2/\delta$. Take $a=0.001$ and $\delta=0.02$, so that $\varpi_R=0.048$. Raising the entry probability to $0.002$, so that one more neutral item in a thousand is followed by a borderline one, raises long-run exposure to $\varpi_R=0.091$. The bound in part (b), with $\|P_{\mathrm{alg}}-P_{\mathrm{ref}}\|_\infty=0.002$, is $0.048$, against an exact change of $0.043$. The derivative is $(1-\varpi_R)^2/\delta=45.4$, and the secant over this change is $43.3$. The magnification is the dwell time $1/\delta=50$ items per visit to $R$, discounted by $(1-\varpi_R)^2=0.91$.
\end{example}

Part (e) shows that sessions cap the effect. In a session of $T$ items, a nudge changes expected exposure by at most $(T-1)/4$ times its size, however slowly the chain mixes. In the long run the change can reach $\kappa_{\mathcal A}$ times its size, and $\kappa_{\mathcal A}$ can be as large as $K/2$. An audit measures $e_T$, not $\varpi$. When sessions are short relative to the time the reference chain needs to reach the antisemitic states, the additive effect measured on audited sessions can understate the long-run effect by a large factor. The amplification ratio, by contrast, can lie on either side of its long-run value. In panel B of \cref{tab:markov-illus}, under the sticky reference, the nudge adds $3.5$ percentage points to $e_{20}$ and $5.6$ to long-run exposure, while the ratio is $5.1$ for $20$-item sessions and $3.4$ in the long run. Extrapolating from the audited horizon to others is the job of the model, which is why its estimation matters.

\subsection{Estimation from audit sessions}\label{sec:markov-estimation}

Within one arm of an audit, fix the policy and the persona and observe $n$ independent sessions $(X_{s1},\dots,X_{sT})$, $s=1,\dots,n$. Let $n_x$ be the number of sessions that start in $x$, let $N_{xy}=\sum_s\sum_{t<T}\ind{X_{st}=x,X_{s,t+1}=y}$ count transitions, and let $N_x=\sum_yN_{xy}$. Let $\nu_x=\sum_{t=1}^{T-1}(\alpha^{\top}P^{t-1})_x$ be the expected number of departures from $x$ per session, and $\Sigma_x=\operatorname{diag}(P_{x\cdot})-P_{x\cdot}^{\top}P_{x\cdot}$, where $P_{x\cdot}$ is the $x$th row of $P$ as a row vector. For a function $\theta(\alpha,P)$, let $\bm c=\partial\theta/\partial\alpha$ and let $G_{x\cdot}=\partial\theta/\partial P_{x\cdot}$, a gradient determined up to a constant in each row.

\begin{proposition}[Estimation]\label{prop:markov-mle}
Let \cref{asm:markov} hold within the arm, and let $\nu_x>0$ for every $x$.
\begin{enumerate}[(a)]
  \item The maximum-likelihood estimators are $\hat\alpha_x=n_x/n$ and $\hat P_{xy}=N_{xy}/N_x$.
  \item As $n\to\infty$ with $T$ fixed, $\sqrt n(\hat\alpha-\alpha)$ and the rows $\sqrt n(\hat P_{x\cdot}-P_{x\cdot})$ are asymptotically independent and normal, with covariance matrices $\operatorname{diag}(\alpha)-\alpha\alpha^{\top}$ and $\Sigma_x/\nu_x$.
  \item If $\theta$ is differentiable at $(\alpha,P)$, then $\sqrt n\{\theta(\hat\alpha,\hat P)-\theta(\alpha,P)\}$ converges in distribution to $\mathrm N(0,\sigma^2_\theta)$, with
  \[
    \sigma^2_\theta=\Var_{x\sim\alpha}(c_x)+\sum_x\frac{1}{\nu_x}\Var_{y\sim P(x,\cdot)}(G_{xy}).
  \]
  \item For long-run exposure, $\theta=\varpi(\mathcal A)$, one may take $\bm c=\bm0$ and $G_{xy}=\varpi_x\,g_{\mathcal A}(y)$, so that $\sigma^2_\theta=\sum_x(\varpi_x^2/\nu_x)\Var_{y\sim P(x,\cdot)}\{g_{\mathcal A}(y)\}$. For session exposure, $\theta=e_T$, one may take $c_x=T^{-1}\sum_{t=0}^{T-1}(P^{t}\bm1_{\mathcal A})_x$ and $G_{xy}=T^{-1}\sum_{t=1}^{T-1}\sum_{s=0}^{t-1}(\alpha^{\top}P^{s})_x(P^{t-1-s}\bm1_{\mathcal A})_y$.
\end{enumerate}
\end{proposition}

Parts (a) and (b) are the classical theory of inference for Markov chains. \Citet{anderson1957} treat exactly this design, many independent sequences of fixed length, and \citet{billingsley1961} surveys the case of one long chain. The likelihood factorises into one multinomial kernel per row, and each row behaves asymptotically like a multinomial sample of size $n\nu_x$. Part (c) is the delta method \citep[Ch.~3]{vandervaart1998}, and part (d) supplies the gradients, the first from \cref{thm:perturb}(a). The two arms of an audit are independent, so the log amplification ratio $\log\{\hat\varpi_{\mathrm{alg}}(\mathcal A)/\hat\varpi_{\mathrm{ref}}(\mathcal A)\}$ has asymptotic variance $\sigma^2_{\mathrm{alg}}/\{n_{\mathrm{alg}}\varpi_{\mathrm{alg}}(\mathcal A)^2\}+\sigma^2_{\mathrm{ref}}/\{n_{\mathrm{ref}}\varpi_{\mathrm{ref}}(\mathcal A)^2\}$, with the $\sigma^2$ of part (d) for each arm. Two observations guide the use of these estimators.

First, the model buys extrapolation more than precision. If the data were one long stationary chain, $\nu_x$ would be proportional to $\varpi_x$, and part (d) would reduce to $\sigma^2_{\mathcal A}$ of \cref{prop:markov}(a): the plug-in estimator of $\varpi(\mathcal A)$ would be exactly as precise as the empirical frequency of antisemitic items. For $e_T$ at the audited horizon, the plug-in $e_T(\hat\alpha,\hat P)$ is asymptotically efficient if \cref{asm:markov} holds, and so at least as precise as the session average $n^{-1}\sum_sE_s$. The gain is often modest, and the plug-in is biased if the model is wrong. The session average should therefore be the primary estimator at the audited horizon, and the chain should be used for what the session average cannot do: the long run, other horizons, other starting distributions, and the sensitivity analysis of \cref{thm:perturb}.

Second, the states are observed through a classifier. If classification errors are independent across items given the true states, with confusion matrix $C$ ($C_{xy}$ is the probability that an item in true state $x$ is classified as $y$), the observed label sequence is a hidden Markov chain, not a Markov chain, and $\hat P$ computed from labels is biased. Its expected transition counts per session are $C^{\top}\operatorname{diag}(\bm\nu)PC$. When $C$ is invertible, $\operatorname{diag}(\bm\nu)P=(C^{\top})^{-1}\,\E(\text{observed counts}/n)\,C^{-1}$, which yields a corrected estimator of $P$ by row normalisation. This is the transition-matrix analogue of the prevalence correction of \citet{rogangladen1978}, and \cref{sec:measurement} treats the estimation of $C$ and differential error.

\subsection{An illustrative example}\label{sec:markov-example}

\Cref{tab:markov-illus} applies \cref{prop:markov,thm:perturb} to a three-state chain with states $N$ (neutral), $B$ (borderline) and $A$ (antisemitic), $\mathcal A=\{A\}$, and sessions that start in $N$. \textbf{All matrices are invented for exposition.} They describe no platform, and the example makes no empirical claim. Rows and columns are ordered $(N,B,A)$:
\[
  P_{\mathrm{ref}}=\begin{pmatrix}.95&.04&.01\\.50&.40&.10\\.40&.30&.30\end{pmatrix},\quad
  P_{\mathrm{alg}}=\begin{pmatrix}.92&.06&.02\\.30&.55&.15\\.20&.35&.45\end{pmatrix},
\]
\[
  P_{\mathrm{rh}}=\begin{pmatrix}.92&.06&.02\\0&.70&.30\\0&.45&.55\end{pmatrix},\quad
  P_{\mathrm{stk}}=\begin{pmatrix}.994&.005&.001\\.05&.89&.06\\.05&.20&.75\end{pmatrix}.
\]
$P_{\mathrm{alg}}$ shifts at most $0.2$ of each row's probability, $\|P_{\mathrm{alg}}-P_{\mathrm{ref}}\|_\infty=0.4$. $P_{\mathrm{rh}}$ is a rabbit-hole variant in which borderline and antisemitic content never lead back to neutral content, so that $\{B,A\}$ is a closed class. $P_{\mathrm{stk}}$ is a sticky reference with almost the same long-run exposure as $P_{\mathrm{ref}}$, but in which neutral content rarely leads elsewhere and the borderline region is hard to leave.

\begin{table}[t]
\centering
\caption{Exposure functionals and sensitivity in a three-state chain (\cref{prop:markov,thm:perturb}). \textbf{Illustrative, invented matrices; not data.} Sessions start in $N$; $\mathcal A=\{A\}$. Panel A: the ratio column is $\mathrm{alg}/\mathrm{ref}$. Panel B: a nudge redirects a fraction $\varepsilon=0.01$ of recommendations to $A$ (\cref{thm:perturb}(c)). All values are exact computations from the matrices in the text.}\label{tab:markov-illus}
\small
\begin{tabular}{@{}l rrr r@{}}
\toprule
\multicolumn{5}{@{}l}{\emph{Panel A. Exposure functionals}}\\
 & $P_{\mathrm{ref}}$ & $P_{\mathrm{alg}}$ & ratio & $P_{\mathrm{rh}}$\\
\midrule
Long-run exposure $\varpi(\mathcal A)$ & 0.0232 & 0.0711 & 3.07 & 0.400\\
Session exposure $e_{20}$ & 0.0207 & 0.0599 & 2.89 & 0.189\\
Session exposure $e_{100}$ & 0.0227 & 0.0689 & 3.03 & 0.348\\
$\Pr(\text{at least one } A \text{ item in } 20)$ & 0.254 & 0.489 & 1.93 & 0.739\\
Items until first $A$ item, from $N$, $M_{N\mathcal A}$ & 64.0 & 28.3 & 0.44 & 15.0\\
Asymptotic variance $\sigma^2_{\mathcal A}$ & 0.044 & 0.189 & & 0.400\\
Kemeny's constant $K$ & 3.13 & 4.27 & & \\
Sensitivity coefficient $\kappa_{\mathcal A}$ & 0.742 & 1.008 & & \\
\midrule
\multicolumn{5}{@{}l}{\emph{Panel B. A nudge of size $\varepsilon=0.01$ toward $A$}}\\
 & $P_{\mathrm{ref}}$ & $P_{\mathrm{stk}}$ & & \\
\midrule
Kemeny's constant $K$ & 3.13 & 21.1 & & \\
Sensitivity coefficient $\kappa_{\mathcal A}$ & 0.742 & 3.31 & & \\
Mean items to $A$ from equilibrium, $\E_\varpi[\tau_A]$ & 61.9 & 261.8 & & \\
Memoryless benchmark $(1-\varpi_A)/\varpi_A$ & 42.1 & 41.3 & & \\
Long-run exposure $\varpi(\mathcal A)$: before nudge & 0.0232 & 0.0236 & & \\
\quad after nudge, exact & 0.0375 & 0.0796 & & \\
\quad after nudge, first order $\varpi_A(1+\varepsilon\,\E_\varpi[\tau_A])$ & 0.0376 & 0.0855 & & \\
\quad after nudge, memoryless $\varpi_A+\varepsilon(1-\varpi_A)$ & 0.0330 & 0.0334 & & \\
Session exposure $e_{20}$: before nudge & 0.0207 & 0.0086 & & \\
\quad after nudge & 0.0339 & 0.0439 & & \\
Session exposure $e_{100}$: before nudge & 0.0227 & 0.0192 & & \\
\quad after nudge & 0.0368 & 0.0702 & & \\
\bottomrule
\end{tabular}
\end{table}

Panel A shows how the algorithm's moderate changes compound. No row of $P_{\mathrm{alg}}$ moves more than $0.2$ of probability, and the largest change in the probability of going to $A$ is $0.15$. Yet long-run exposure triples, from $2.3\%$ to $7.1\%$. The exact identity of \cref{thm:perturb}(a) reproduces $\varpi_{\mathrm{alg}}(\mathcal A)=0.0711$ from $P_{\mathrm{ref}}$ and the change alone. The first-order prediction is $0.030$ against an exact change of $0.048$, because the change is not small. The global bound $\kappa_{\mathcal A}\|P_{\mathrm{alg}}-P_{\mathrm{ref}}\|_\infty=0.297$ holds with room to spare, and the Kemeny bound is $0.626$. The bound is nearly attained in \cref{ex:mk-two-region} but is loose here, because the change is spread over directions that do not all point at the exposure potential's extremes. The ratio of session exposures at $T=20$ is $2.89$, below the long-run ratio of $3.07$, because sessions start in neutral territory (\cref{prop:markov}(b)). Reach, the probability of at least one antisemitic item, rises by less, from $25\%$ to $49\%$. In the rabbit-hole variant every session ends in $\{B,A\}$, after $12.5$ items on average, and long-run exposure is $40\%$ whatever the entry probability. A $20$-item session has not yet reached that level: $e_{20}=0.189$.

Panel B isolates stickiness. The two references have nearly the same long-run exposure, about $2.3\%$, and the same memoryless benchmark. Redirecting one recommendation in a hundred to $A$ raises long-run exposure to $3.7\%$ under $P_{\mathrm{ref}}$, close to the memoryless $3.3\%$. Under the sticky reference it raises it to $8.0\%$, an increase of $5.6$ percentage points from a nudge of one percentage point. This is \cref{thm:perturb}(c) at work: the mean time to reach $A$ from equilibrium is $262$ items under $P_{\mathrm{stk}}$ against $62$ under $P_{\mathrm{ref}}$. By \cref{thm:perturb}(d), with $S=\{N\}$ and $q(S)=0.006$, it is at least $149$. The first-order prediction overshoots slightly, $8.5\%$, because a nudge of this size is no longer small for so sensitive a chain. Finite sessions show the same ordering. Under the sticky reference, $20$-item sessions see five times as much antisemitic content after the nudge as before.
